% Generated draft: requires mathematical and source review.
\documentclass[11pt]{article}
\usepackage[margin=25mm]{geometry}
\usepackage{fontspec}
\setmainfont{Latin Modern Roman}
\usepackage{amsmath,amssymb,mathtools,unicode-math}
\setmathfont{Latin Modern Math}
\usepackage{enumitem,needspace,xurl,hyperref}
\setlength{\parindent}{0pt}
\setlength{\parskip}{5pt}
\setlength{\emergencystretch}{3em}
\newcommand{\N}{\mathbb{N}}
\newcommand{\Z}{\mathbb{Z}}
\newcommand{\Q}{\mathbb{Q}}
\newcommand{\R}{\mathbb{R}}
\newcommand{\C}{\mathbb{C}}
\newcommand{\F}{\mathbb{F}}
\newcommand{\A}{\mathbb{A}}
\newcommand{\PP}{\mathbb{P}}
\newcommand{\E}{\mathbb{E}}
\newcommand{\eps}{\varepsilon}
\newcommand{\dd}{\,\mathrm{d}}
\begin{document}
\section*{UnsolvedMath ID: 30006370}
\textbf{Source problem code:} catalogue entry
\subsection{Definitions and mathematical statement}
The supplied record asks about rational points near a submanifold \(M\), described in its background assessment as a bounded immersed submanifold. It does not specify the ambient space, the dimension or regularity of \(M\), the height convention, or the meaning and size of the neighborhood.

A conditional formulation of the counting question is as follows. Once an ambient space \(X\), its designated set of rational points \(\mathcal R_X\), a height function \(h\), and a prescribed neighborhood \(U\subseteq X\) of \(M\) have been specified, define
\[
N_{M,U,h}(H)
:=\#\bigl\{P\in\mathcal R_X\cap U:h(P)\le H\bigr\},
\]
where \(H\) is the height bound and \(\#\) denotes cardinality, without an assertion here that it is finite. The question is to determine how many such points there are, with dependence on the height bound, the prescribed small neighborhood, and the submanifold.

This is a schematic counting problem rather than a specified conjectural bound. The original wording describes height as a bound on denominator size, but does not define denominators or their normalization. No particular ambient rational structure, metric, neighborhood convention, dependence of neighborhood size on \(H\), curvature hypothesis, nondegeneracy condition, or desired estimate is supplied. These choices remain unresolved; the display defines the counting quantity only after they have been fixed and does not assert a universal estimate.
\subsection{Short English statement}
Count rational points of height at most \(H\) in a prescribed small neighborhood of a bounded immersed submanifold \(M\). The supplied record leaves the ambient space, height, neighborhood, geometric hypotheses, and target bound unspecified, so it supports only a conditional counting formulation.
\subsection{Sources}
\begin{itemize}
\item Supplied local UnsolvedMath registry record, permanent ID \(30006370\), title \emph{Rational Points Near Submanifolds}, collection \emph{Oberwolfach Reports Open Problems}, exact source problem code \texttt{catalogue entry}.
\item Supplied source citation: \emph{Real Analysis, Harmonic Analysis, and Applications (2026)}, Oberwolfach Reports, DOI \texttt{10.4171/owr/2025/32}, pp.\ 1707--1752. \url{https://doi.org/10.4171/owr/2025/32}. The original statement appends the talk title \emph{Counting, Curvature and Convex Duality} and the name Rajula Srivastava; the separate proposer field is empty.
\item Supplied literature source labeled ``primary workshop text'': \url{https://oa.tib.eu/renate/server/api/core/bitstreams/fdbbbd0a-d4c1-475f-aaa8-2151151a7f4b/content}.
\item Literature assessment dated exactly \texttt{2026-08-22}: the supplied curation labels the question open and describes it as concerning rational points of bounded height near a bounded immersed submanifold, with no general proved assertion or hypotheses extracted. This is supplied dated curation, not an independently verified literature assessment. The background expressly requests independent verification of time-sensitive status and cited claims.
\item The supplied citation labels the report as 2026, while the DOI contains 2025 and the proposed-year field and workshop-year tag specify 2025. The supplied record does not resolve this dating discrepancy.
\end{itemize}
\end{document}
